% ════════════════════════════════════════════
% 基本假设（03.模型假设.tex）
% ════════════════════════════════════════════
\subsection{基本假设}

为明确分析边界，本报告做出以下假设。

\textbf{假设1：} 同一主页的 500 条帖子在给定协变量下相互独立，单条帖子的曝光不显著挤占其他帖子的推荐流量，因而可作为独立样本做统计推断。

\textbf{假设2：} 付费标记的 1 条缺失为完全随机缺失，众数填补不改变参数方向与量级，缺失处理方案的差异可忽略。

\textbf{假设3：} log1p 曝光量的条件均值对解释变量呈可加线性关系，月份固定效应能够吸收粉丝规模随时间累积的趋势。

\textbf{假设4：} 在控制内容类型、主题与发布时机后，付费标记与内容类型同随机扰动不相关，反事实预测用于情景比较而非严格的因果识别。

上述假设不改变数据本身，只界定结论的适用范围。若独立性或完全随机缺失不成立，区间估计会偏窄；若线性设定不成立，系数应理解为局部平均效应而非全局因果效应。报告在相应位置说明了这些限制，并用六种设定与缺失处理对照检验结论对假设的敏感程度。
